\section{Computational methods}\label{app:code}
All computations were done with the author's Python code (Koornwinder polynomials and the polynomial representations of the DAHA of types $D_N$, $C_N$ and $\GL_N$; rational function arithmetic with FLINT). Every Result in the text is checked by an automated test; the checks include the reproduction of the values of \cite{Marino09}, agreement with Kauffman polynomials at $t = q$, agreement with direct fixed-$N$ computations, agreement with the invariants of \cite{KN17}, and the symmetry $R \leftrightarrow R^T$. Two identifications used as input were verified numerically as checks of the implementation: the dual Cauchy identity \eqref{eq:mimachi} for $(n, m) = (1,1), (2,1), (1,2), (2,2)$ (automated tests) and $(2,3), (3,2)$ (direct runs of the same check), and the identification of the $D_N$ Macdonald dimension $\dim_R$ of \cite{AgSc12} with the principal specialization of the Koornwinder polynomial for the representations $(1), (2), (1,1), (3), (2,1), (1,1,1)$ at $N = 2, 3$ (against the evaluation formulas of \cite{vD96,Sahi99} and against polynomials constructed by the $q$-difference operator).

\subsection{DAHA conventions and checks}\label{app:daha-conv}
In \eqref{eq:JD} the conventions are: $(T_i - t^{1/2})(T_i + t^{-1/2}) = 0$; since the vector representation is minuscule, $Y_{\varepsilon_1} = \pi_1 T_{u_1}$ with $\pi_1(X_b) = q^{(\omega_{1*}, b)}X_{u_1^{-1}b}$ ($\omega_{1*} = \varepsilon_1$ for $D_N$, $\omega_{1*} \equiv -\varepsilon_N$ for $\GL_N$); $\tau_-(X) = q^{1/2}YX$, $\tau_+(Y) = q^{-1/2}XY$, and $T(2,2m+1)$ corresponds to $\gamma = \tau_+^m\tau_-^2$. To speed up the computation we use the anti-involution $\varphi$ exchanging $X$ and $Y$, symmetrize over orbits, and share intermediate results between representations. The composites $(\alpha;\bar1)$ of the cover are computed by starting the chain of operators on the $X$ side from $\hat\gamma'(p_1(Y))\cdot 1$, where $\hat\gamma' = \varphi\hat\gamma\varphi$ and $p_1(Y)\cdot 1$ is built by symmetrization on the $X^{-1}$ side. For four boxes one needs $|\alpha| = 3$, where the irreducible decomposition $P_\alpha(X)p_1(X^{-1}) = P_{(\alpha;\bar1)}(X) + \sum_\mu c_{\alpha,\mu}P_\mu(X)$ ($\mu = \alpha - \varepsilon_i$) has several lower terms. From $p_1(X^{-1}) = e_{N-1}(X)/e_N(X)$ and the Pieri rule for vertical strips, $c_{\alpha,\mu} = \psi'_{\nu/\alpha}$ with $\nu = \alpha + 1^N - \varepsilon_i$ (the Pieri coefficient $\psi'$ of Chapter VI of \cite{Macdonald95}); only the leg in column $0$ depends on $N$, which gives a closed form in $t^N$, checked against the polynomial decomposition for small $N$ and against the identity of principal specializations for all $N$. For iterated torus knots (\S\ref{sec:cable}) the polynomial from the pattern step is expanded in the basis values $\{\hat\gamma_{\mathrm{comp}}(H_\nu)\}$ of the companion step, which are computed with coefficients kept as Laurent polynomials without cancellation.
% [D1][D2][D3]
\begin{result}[checks]\label{cr:daha}
The following were verified.
\begin{enumerate}
\item For $D_N$, $\GL_N$ ($N = 2, 3, 4$) the DAHA relations (quadratic, braid, cross, commutativity of $Y$) hold, and the $D_N$ Macdonald polynomials are eigenfunctions of the Macdonald operator. For $C_2$, $C_3$ the DAHA relations hold as well.
\item $D_N$: $Z = T_R^{-F}\,\mirror(W_R)$ with $W_R$ the fixed-$N$ invariant in the uniform normalization; checked for $N = 2, 3$, $R \in \{(1), (1,1), (2)\}$, $m = 1, 2$. For $N = 4, 5$ the fundamental representation coincides exactly with the mirror of (6.15) of \cite{AgSc12}.
\item $\GL_N$: \eqref{eq:Mb} with $\cK = m+1$. For $m = 1$ the rule was read off from $(1), (2), (1,1)$ and predicted correctly for $(3), (2,1), (1,1,1)$ and the adjoint ($\U(3)$, $\U(5)$); for $m = 2$ it holds for $(1), (2), (1,1)$ at $N = 2, 3, 4$ (nine cases) and for the adjoint. For composites with $|\alpha| = 3$ the values of $(3;\bar1)$ in $\U(2)$ and $(2,1;\bar1)$, $(1,1,1;\bar1)$ in $\U(4)$ agree exactly with those of $(4)$, $(3,2,1)$, $(2,2,2)$ ($m = 1$).
\item The fast and the reference implementations agree coefficient by coefficient on all computed values.
\end{enumerate}
\end{result}
The heaviest computations are the three-box $T(2,9)$ on the $D_{14}$ DAHA (90 minutes and 16\,GB per item) and one basis value of the $(2,13)$-cable on $D_{15}$ (about 2 hours and 19\,GB).

\subsection{Interpolation in \texorpdfstring{$N$}{N}}\label{app:interp}
The reduced invariant $\Rred_R(N) = T_R^{-F}\,\mirror(Z(N))/K_R(t^\rho)$ (with $Z$ the unnormalized value of \S\ref{sec:daha}) is a Laurent polynomial in $U = t^N$ whose coefficients are Laurent polynomials in $q^{1/2}, t^{1/2}$. Its minimal and maximal degrees in $t^{1/2}$, as functions of $N$, become linear, $2Nj_{\min} + \mathrm{const}$ and $2Nj_{\max} + \mathrm{const}$, once $N$ is large enough; the slopes $j_{\min}, j_{\max}$ are read off from consecutive $N$ and give the range of powers $U^j$, $j_{\min} \le j \le j_{\max}$. With $n = j_{\max} - j_{\min} + 1$ unknown coefficients, Newton's divided differences on $n$ nodes $N$ determine them, and the result is compared coefficient by coefficient with the direct computation at the remaining $N$. When $n$ is large (two boxes of $T(2,7)$ and $T(2,9)$ and their adjoints, three boxes of $T(2,5)$ and beyond, four boxes, composites with $m \ge 2$, two boxes of the cables), the unknowns are restricted degree by degree in $q^{1/2}$: for each degree $\alpha$ of $q^{1/2}$ and each power $U^j$, the coefficient is assumed to have its degrees in $t^{1/2}$ within the range of degrees of the $q^{\alpha/2}$-part of the computed value at a reference $N_0$ (a ``band''); the resulting linear system over all nodes is solved and accepted only if it has a unique solution, and the solution is then compared with the values at all $N$ not used. For the $(1,1)$ invariant of $\Cab{13}$ the band was chosen separately for each degree $\alpha$, with the smallest $N_0$ for which the system is consistent. The verified ranges are in Tables \ref{tab:interp} and \ref{tab:cableinterp}.
% [D4]
\begin{result}\label{cr:interp}
In the ranges of Tables \ref{tab:interp} and \ref{tab:cableinterp} the values at the $N$ not used in the interpolation were predicted coefficient by coefficient. At $t = q$ the three-box $(3)$ agrees exactly with the mirror of the $(3)$-colored Kauffman polynomial of \cite{NRZ13}, and the four-box $(4)$ and $(1,1,1,1)$ of the trefoil agree with the values of \cite{NRZ13}.
\end{result}
\begin{table}[htbp]
\centering
\caption{Interpolation in $N$ of the reduced invariants of torus knots. ``Powers'' are the powers of $t^N$ on the $\SO$ side and of $\lambda = t^{N_A/2}$ for the cover; ``nodes'' are the values of $N$ (or $N_A$) used in the interpolation, and ``checked at'' the values at which the result was compared with the direct computation. In rows with several representations the nodes are listed in the order of the representations. Rows below the rule use the restricted-range method.}\label{tab:interp}
\scriptsize
\setlength{\tabcolsep}{2pt}
\resizebox{\textwidth}{!}{%
\begin{tabular}{llll}
\toprule
knot, representation & powers & nodes & checked at \\
\midrule
trefoil, $(1)$ & $-5$ to $-2$ & $N = 2$--5 & $N = 6, 7$; (6.15) of \cite{AgSc12} \\
trefoil, $(2), (1,1)$ & $-10$ to $-4$ & $N = 2$--8 & $N = 9$ (and $N = 2$ from $N = 3$--9) \\
trefoil, $(3), (2,1), (1,1,1)$ & $-15$ to $-6$ & $N = 4$--13 & $N = 2, 3$ ($N = 3$ for $(1,1,1)$) \\
$T(2,5)$, $(1)$ & $-9$ to $-4$ & $N = 9$--14 & $N = 2$--8 \\
$T(2,5)$, $(2), (1,1)$ & $-18$ to $-8$ & $N = 4$--14 & $N = 2, 3$ \\
cover adjoint $(1;\bar1)$, trefoil & $-10$ to $-4$ & 7 of $N_A = 2$--11 & the other 3 \\
cover adjoint $(1;\bar1)$, $T(2,5)$ & $-18$ to $-8$ & $N_A = 14$--24 & $N_A = 2$--13 \\
cover $(2;\bar1), (1,1;\bar1)$, trefoil & $-14$ to $-6$ & $N_A = 12$--20 & $N_A = 2$ (3 for $(1,1;\bar1)$)--11 \\
\midrule
$T(2,7)$, $(2), (1,1)$ & $-26$ to $-12$ & $N = 3$--8, 4--9 & $2, 9$--20; $2, 3, 10$--20 \\
$T(2,9)$, $(2), (1,1)$ & $-34$ to $-16$ & $N = 3$--8, 4--10 & $2, 9$--26; $2, 3, 11$--26 \\
cover adjoint $(1;\bar1)$, $T(2,7)$ & $-26$ to $-12$ & $N_A = 3$--8 & $N_A = 2, 9$--26 \\
cover adjoint $(1;\bar1)$, $T(2,9)$ & $-34$ to $-16$ & $N_A = 3$--8 & $N_A = 2, 9$--34 \\
$T(2,5)$, $(3), (2,1), (1,1,1)$ & $-27$ to $-12$ & $N = 4$--10, 3--10, 5--11 & remaining $N \le 14$ \\
$T(2,7)$, $(3), (2,1), (1,1,1)$ & $-39$ to $-18$ & $N = 2$--9, 3--12, 5--12 & remaining $N \le 13$ \\
$T(2,9)$, $(3), (2,1), (1,1,1)$ & $-51$ to $-24$ & $N = 2$--10, 3--13, 5--13 & remaining $N \le 14$ \\
trefoil, four boxes (five $R$) & $-20$ to $-8$ & 7--9 of $N \le 12$ & remaining $N \le 12$ \\
cover $(2;\bar1), (1,1;\bar1)$, $m = 2, 3, 4$ & from $-12$ to between $-24$ and $-40$ & $N_A \le 18$--20 & remaining $N_A$ \\
cover $(3;\bar1), (2,1;\bar1), (1,1,1;\bar1)$, trefoil & $-18$ to $-8$ & $N_A = 3$--8, 5--11, 4--12 & remaining $N_A \le 14$ \\
\bottomrule
\end{tabular}}
\end{table}
\begin{table}[htbp]
\centering
\caption{Interpolation in $N$ for the cables (same conventions as Table \ref{tab:interp}; all rows use the restricted-range method except the fundamental representation).}\label{tab:cableinterp}
\footnotesize
\begin{tabular}{llll}
\toprule
object & powers & nodes & checked at \\
\midrule
$\Cab{13}$, $(1)$ & $-25$ to $-16$ & $N = 3$--8 & $N = 2, 9$--16 \\
$\Cab{15}$, $(1)$ & $-29$ to $-18$ & $N = 4$--10 & $N = 2, 3, 11$--16 \\
$\Cab{13}$, $(2)$ & $-50$ to $-32$ & $N = 3$--11 & $N = 2, 12$--15 \\
$\Cab{13}$, $(1,1)$ & $-50$ to $-32$ & $N = 4$--14 & $N = 2, 3, 15$ \\
$\Cab{15}$, $(2)$ & $-58$ to $-36$ & $N = 2$--9 & $N = 10$ \\
cover adjoint $(1;\bar1)$, $\Cab{13}$ & $-50$ to $-32$ & $N_A = 5$--14 & $N_A = 2, 3, 4$ \\
cover adjoint $(1;\bar1)$, $\Cab{15}$ & $-58$ to $-36$ & $N_A = 5$--14 & $N_A = 2, 3, 4$ \\
\bottomrule
\end{tabular}
\end{table}

\subsection{One-probe computation for the Hopf link}\label{app:virtual}
For the Hopf link with one probe on each component we use the probe basis $K_{(n)}(y; t, q; 1, -1, q^{1/2}, -q^{1/2})$ (monic Askey--Wilson polynomials).
\begin{conjecture}\label{conj:c2j}
The coefficient of $y^{n-2j}$ in $K_{(n)}(y; t, q; 1, -1, q^{1/2}, -q^{1/2})$ is, as a rational function of $t^n$,
\begin{equation}\label{eq:c2j}
c_{2j}(n) = t^j\, \frac{(q/t; t^2)_j}{(t^2; t^2)_j}\cdot\frac{(1 - t^n)(1 + t^{n-2j})\prod_{i=1}^{j-1}(1 - t^{2n-2i})}{\prod_{i=1}^{j}(1 - qt^{2n-2i-1})} .
\end{equation}
\end{conjecture}
% [S4]
Equation \eqref{eq:c2j} agrees with the recursion for the Askey--Wilson coefficients at four rational points for $j \le 5$, $n \le 30$. With it, the virtual normalization $Z^{\mathrm{virt}} = y_1^N y_2^N Z$ is a polynomial in $y_1, y_2$ whose coefficient of $y_1^a y_2^b$ is a finite sum independent of $N$, so $U$ can be kept as a formal variable in this one-probe computation (full agreement with the direct computation at $U = t^N$, $N = 3, \dots, 8$, for $a + b \le 8$). The multi-probe computations of \S\ref{sec:kernel} do not use \eqref{eq:c2j}.

\subsection{Integer programs for \texorpdfstring{$E$}{E}}\label{app:ilp}
When $E$ is said to be determined by (E1) and (E2), the following problem was solved. Let $\Phifix = \sum_\omega c_\omega\,\omega$ be the index of the fixed sector (a Laurent polynomial in $q^{1/2}, t^{1/2}, \lambda$) and write $E = \sum_\omega e_\omega\,\omega$ with unknown integers $e_\omega$ over the monomials of $\Phifix$ (or of a stated larger support). The constraints are: (i) $e_\omega \equiv c_\omega \pmod 2$ (integrality of $f$, Remark in \S\ref{sec:EC}); (ii) $|e_\omega| \le |c_\omega|$ (Hypothesis \ref{hyp:QP}, used for the Hopf link; for knots this bound is replaced by the explicit ans\"atze of Appendix \ref{app:attempts}); (iii) the linear equations (E1), $\sum_\omega e_\omega\,\omega|_{t=-1} = [2(1+t)g]_{t=-1}$ coefficient by coefficient; (iv) when (E2) is imposed, the linear equations $\sum_\omega e_\omega\,\omega|_{t=q} = -(1+q)\psi^2 g^{\mathrm{odd}}$ coefficient by coefficient. Feasibility and uniqueness are decided by solving the integer linear program with an exact solver; ``unique'' means that the solution set has one element, and ``$k$ of $n$ monomials free'' means that the affine solution set has dimension $k$ within the $n$ unknowns. The minimal number of exchanged pairs quoted in Appendix \ref{app:n2} is the minimum of $\tfrac12\sum_\omega(|c_\omega| - |e_\omega|)$ over the solution set.
